\addcontentsline{toc}{subsection}{Inciso e}

\subsection*{e.}

\textbf{Durante el desarrollo de un nuevo sistema empresarial surge una discusión entre quienes consideran que la base de datos debe diseñarse pensando únicamente en las necesidades actuales y quienes proponen anticipar futuros requerimientos. Analiza esta situación relacionándola con el ciclo de vida de una base de datos y argumenta cómo las decisiones tomadas durante las etapas iniciales pueden influir en el mantenimiento, evolución y costo del sistema a largo plazo.}\\ \\
El argumento para diseñar el sistema pensando únicamente en los requerimientos actuales sería ahorrar tiempo en la fase de diseño, para poder implementar la base de datos lo más pronto posible.\\ \\
Sin embargo, las decisiones tomadas durante la fase de diseño determinarán qué tan costoso será modificar el sistema después. Anticiparse a las necesidades predecibles futuras permite que durante el mantenimiento y las futuras intenciones de escalar la base de datos, los cambios realizados no afecten los datos históricos ni los sistemas que ya funcionan sobre el diseño original, evitando costos elevados en migraciones de datos y reescritura de aplicaciones, además de posibles interrupciones en el funcionamiento del sistema.